\begin{abstract}
Highway traffic jams are a congestion-control failure in a transportation system: small local disturbances push traffic past a critical density, after which stop-and-go waves persist long after their cause has disappeared. Field experiments show that a handful of controlled vehicles can damp these waves, and a recent self-regulating-cars system improves throughput by 5\% and cuts average delay by 13\% on a real highway network in high-fidelity simulation, yet deployed designs still place purpose-built vehicles, roadside infrastructure, or a cloud coordinator inside the control loop. We argue the remaining obstacle is deployment, and ask two questions: how little must change in an ordinary car for it to self-regulate, and how few self-regulating cars does a road network need? We describe a vehicle that senses traffic with a dashcam-grade camera and a GPS or OBD-II speed feed, maintains an ego-centric digital twin on roughly \$250 of embedded compute, and independently computes safety-filtered speed and headway advisories, using a low-rate aggregate vehicle-to-vehicle exchange to extend its view. Seen this way, a self-regulating fleet is a communication network laid over the road network: edge nodes deciding locally, a low-rate substrate among them, and a cloud confined to training and updates, so what to compute where, what to share, and how much staleness to tolerate are networking questions. A working prototype grounds the design, running a trained policy end to end in 20--25\,ms. Minimal fleets complete the argument: on bottleneck-fed corridors like the Holland Tunnel approaches, about 25 vehicles, the ambient presence of one commercial fleet, can regulate the inflow that decides whether the bottleneck jams. We close with the networking and systems questions this opens.
\end{abstract}
